\printedpage{(non numérotée)}
% head: (none; title page)

\begin{center}
{\Large\textsf{THE DERIVATORS}}

\bigskip

\textbf{ALEXANDRE GROTHENDIECK}

\bigskip\bigskip

\textbf{Chapter XIX}

\medskip

\textbf{Models (4)}
\end{center}

\newpage

\printedpage{(non numérotée)}
% head: (none; editorial note)

This text was deciphered and transcribed in \LaTeXe\ by M. Künzer. It was edited by M. Künzer, J. Malgoire, and G. Maltsiniotis. The transcription is as faithful as possible to the manuscript. Only a few obvious minor corrections have been made. For the rare additions or comments of the editors, as well as for the original numbering of the manuscript pages, typewriter characters \texttt{[typewriter]} in square brackets are used.

This edition is provisional. Remarks, comments and corrections are welcome. Send a message to:

\textbf{[e-mail address removed]}

\hfill G. Maltsiniotis\hspace*{2cm}

\newpage

\printedpage{1}
% head: A. Grothendieck : Dérivateurs, Ch. XIX, Édition des Universités de Montpellier II et Paris VII
% foot: Version 11 février 2019

\msp{page 1}

\section*{1\quad Reminders on Quillen couples (special Quillen couples)}

I am going to take up again the notations and terminology for Quillen's LLP \texttt{[left lifting property]} and RLP \texttt{[right lifting property]} relations. I write (for a category $\mathcal{M}$, and $\Phi, \Psi \subset \mathrm{Fl}(\mathcal{M})$)
\[
\Phi \Qo \Psi \qquad \iff \qquad (\Phi \subset \Psi^{*}) \qquad \iff \qquad (\Psi \subset \Phi_{*})\ ,
\tag*{(1.1)}
\]
i.e. $\Phi$ has the LLP with respect to $\Psi$, or again $\Psi$ the RLP with respect to $\Phi$. I also say that $(\Phi, \Psi)$ is a \textit{couple} which is Q-\textit{orthogonal.} This notion is not symmetric, but I sometimes allow myself (exceptionally) to write
\[
\Psi \Qorev \Phi \qquad \text{as a synonym for} \qquad \Phi \Qo \Psi\ .
\tag*{(1.2)}
\]
$\Phi$ is Q-\textit{orthogonal on the left} to $\Psi$, $\Psi$ is Q-\textit{orthogonal on the right} to $\Phi$, and Q-orthogonality on the left (resp. on the right) is denoted by the half-arrowhead pointing downward (resp. upward). I write
\[
\left\{
\begin{aligned}
&\Phi \quad \Qog \quad \Psi \quad \iff \quad (\Phi = \Psi^{*})\\
&\Phi \quad \Qod \quad \Psi \quad \iff \quad (\Psi = \Phi_{*})\\
&\Phi \quad \Qogd \quad \Psi \quad \iff \quad (\Phi = \Psi^{*} \ \text{ and } \ \Psi = \Phi_{*})\ .
\end{aligned}
\right.
\tag*{(1.3)}
\]
These three cases of strengthened Q-orthogonality are thus those where we have in addition $\Phi = \Phi^{\tld}$
resp. $\Psi = \Psi^{\rcl}$ resp. $\Phi = \Phi^{\tld}$ and $\Psi = \Psi^{\rcl}$. We say that $\Phi$ is Q-\textit{left closed} if $\Phi = \Phi^{\tld}$, that % reading resolved: DER_XIX_corrections.csv
$\Psi$ is Q-\textit{right closed} if $\Psi = \Psi^{\rcl}$, that the Q-orthogonal couple $(\Phi, \Psi)$ is Q-\textit{left closed} (resp. \textit{right}) if $\Phi$ (resp. $\Psi$) is left Q-closed (resp. right). We say that it is Q-\textit{closed} if the two conditions % reading resolved: DER_XIX_corrections.csv

\msp{page 2}

are satisfied simultaneously. When no confusion is to be feared, I shall sometimes say simply that $\Phi$ is Q-closed, or that $\Psi$ is Q-closed, without specifying whether it is on the left or % [sic] où -> ou, conjunction "or" intended
on the right. But for a Q-orthogonal couple $(\Phi, \Psi)$ (where I prefer to write it $\Phi \Qo \Psi$), Q-closed % [sic] ou -> où, relative adverb intended
always means left and right Q-closed.

When a Q-orthogonal couple $\Phi \Qo \Psi$ has Quillen's factorization property, I indicate it by the vertical bar crossing the horizontal, thus
\[
\Phi \ \Qx{\leftharpoondown}{\rightharpoonup} \ \Psi\ ,
\tag*{(1.4)}
\]
and we define in the same way the meaning of the graphisms
\[
\left\{
\begin{aligned}
&\Phi \quad \Qx{\leftharpoondown}{\rightharpoonup} \quad \Psi\\
&\Phi \quad \Qx{\leftarrow}{\rightharpoonup} \quad \Psi\\
&\Phi \quad \Qx{\leftharpoondown}{\rightarrow} \quad \Psi\\
&\Phi \quad \Qx{\leftarrow}{\rightarrow} \quad \Psi\ .
\end{aligned}
\right.
\tag*{(1.3)} % [sic] (1.3) -> (1.5), number already used above and (1.5) is missing
\]

\newpage

\printedpage{2}
% head: A. Grothendieck : Dérivateurs, Ch. XIX, Édition des Universités de Montpellier II et Paris VII
% foot: Version 11 février 2019

A Q-\textit{orthogonal} Q-\textit{factorizing} couple $\Phi \Qx{\leftharpoondown}{\rightharpoonup} \Psi$ is called a \textit{Quillen precouple.} If it is moreover Q-closed, i.e. if we have $\Phi \Qx{\leftarrow}{\rightarrow} \Psi$, we say that it is a \textit{Quillen couple.}

$\Phi$ is said to be Q-\textit{left saturated} if it contains the isomorphisms and is stable under cobase change, under (small) transfinite composition, and under direct factors. Dually $\Psi$ is said to be Q-\textit{right saturated} if it

\msp{page 3}

contains the isomorphisms, is stable under base change, under transfinite composition, and under direct factors. If $\Phi$ is left Q-saturated, it is stable under (small) direct sums, i.e. if the $f_i : A_i \longrightarrow B_i$ ($i \in I$, $I$ small) are in $\Phi$, $f = \coprod f_i : \coprod A_i \longrightarrow \coprod B_i$ also, and likewise for (small) finite or infinite amalgamated sums (when we work in an $A\backslash\mathcal{M}$ instead of $\mathcal{M}$). We have the dual statement for $\Psi$, when $\Psi$ is right Q-saturated.

In fact, there are weaker notions of Q-saturation than the one I have just given: Q-\textit{simple saturation} (on the left: $\Phi$ contains the isomorphisms, stable under binary composition, under cobase change), Q-\textit{transfinite saturation} (we require in addition stability under transfinite composition), finally Q-\textit{total saturation,} or Q-saturation for short (when we require in addition stability under direct factors). We have the implications
\[
\left\{
\begin{aligned}
&\text{left Q-closed}\\
&\Downarrow\\
&\text{Q-saturated (total)}\\
&\Downarrow\\
&\text{transfinitely Q-saturated}\\
&\Downarrow\\
&\text{simply Q-saturated.}
\end{aligned}
\right.
\tag*{(1.6)}
\]
For $\Phi \subset \mathrm{Fl}(\mathcal{M})$, we have its (left) Q-closure $\Phi^{\tld}$ (which is the smallest left Q-closed subset of $\mathrm{Fl}(\mathcal{M})$ containing $\Phi$) and likewise its (total) left Q-saturation, its transfinite Q-saturation, its simple Q-saturation

\msp{page 4}

with the tautological inclusions
\[
\Phi \ \subset \
\underbrace{\Phi'}_{\substack{\text{Q-saturation}\\ \text{simple on the left}}}
\ \subset \
\underbrace{\Phi''}_{\substack{\text{Q-saturation}\\ \text{transfinite on the left}}}
\ \subset \
\underbrace{\Phi'''}_{\substack{\text{Q-saturation}\\ \text{total on the left}}}
\ \subset \
\underbrace{\Phi^{\tld}}_{\substack{\text{Q-closure}\\ \text{on the left}}}
\ ,
\tag*{(1.7)}
\]
and dually for $\Psi$. Between $\Phi$ and $\Phi'$ there is also a $\Phi'_0$,
\[
\Phi \ \subset \ \Phi'_0 \ \subset \ \Phi'\ ,
\tag*{(1.8)}
\]
which is the Q-\textit{saturation under cobase change} (also containing, by definition, all isomorphisms). To construct all these saturations, we proceed as follows:

\newpage

\printedpage{3}
% head: A. Grothendieck : Dérivateurs, Ch. XIX, Édition des Universités de Montpellier II et Paris VII
% foot: Version 11 février 2019

\begin{enumerate}[leftmargin=*]
\item[1)] $\Phi'_0 \quad = \quad$ set of arrows $i : A \longrightarrow B$ fitting into a cocartesian square
\[
\begin{tikzcd}
A_0 \arrow[r, "i_0"] \arrow[d] & B_0 \arrow[d]\\
A \arrow[r, "i"] & B\ ,
\end{tikzcd}
\]
with $i_0 \in \Phi$ or $i_0$ an identity (in which case we find for $i$ the isomorphisms).

\item[2)] $\Phi' \quad = \quad$ set of multiple composite arrows
\[
i_n \, i_{n-1} \ \cdots \ i_1\ ,
\]
with $i_\alpha \in \Phi'_0$.

\item[3)] $\Phi'' \quad = \quad$ set of transfinite composites of arrows $\in \Phi'_0$.
\end{enumerate}

We feel like putting
\[
\Phi''' \ = \ \Phi^{\prime\prime\natural} \ \overset{\text{def}}{=} \ \text{set of direct factors of elements of } \Phi'',
\]
and it is a fact that $\Phi^{\prime\prime\natural} \subset \Phi'''$ trivially, but it is not clear in general that we have % [sic] par -> pas, typo
\msp{page 5}
equality, for lack of being able to show that $\Phi^{\prime\prime\natural}$ is stable, not only under direct factors, but also under cobase change and under transfinite composition. However, let us recall:

\textbf{Theorem 1.1.} \textit{Suppose $\mathcal{M}$ stable under small $\varinjlim$ , and that its elements are accessible. Suppose moreover that $\Phi$ is small, more generally, that there exists a small subset $\Phi_0 \subset \Phi$ such that $\Phi \subset \Phi_0{}^{\tld}$, hence}
\[
\Psi \ \overset{\text{def}}{=} \ \Phi_{*} \ = \ (\Phi_0)_{*} \qquad \text{\footnotemark}.
\tag*{(1.1.1)} % footnote mark printed as superscript (1)
\]
\footnotetext{\textbf{N.B.} We thus have $\Phi \Qod \Psi$, a fortiori $\Phi_0 \Qod \Psi$.}
\textit{Under these conditions, if $\Phi''_0$ denotes the transfinite left \textup{Q}-saturation of $\Phi_0$, then we have $\Phi''_0 \Qx{\leftharpoondown}{\rightarrow} \Psi$, i.e. every arrow $f$ of $\mathcal{M}$ factors as $f = pi$, with $i \in \Phi''_0$, $p \in \Psi$\,\footnote{There is even \textit{functorial} factorization in $f$, a property which deserves a notation, e.g. $\Phi''_0 \Qxx{\leftharpoondown}{\rightarrow} \Psi$.}. A fortiori, if $\Phi''$ denotes the transfinite left \textup{Q}-saturation of $\Phi$, we have} % footnote mark printed as superscript (2)
\[
\Phi'' \ \Qx{\leftharpoondown}{\rightarrow} \ \Psi\ .
\]

\bigskip

\textbf{Corollary 1.2.} \textit{Under the conditions of the theorem, we have}
\[
\Phi''' \ \text{(total left Q-saturation of $\Phi$)} \ = \
\underbrace{\Phi^{\prime\prime\natural}}_{\substack{\text{set of}\\ \text{direct factors of elements}\\ \text{of $\Phi''$}}}
\ = \ \Phi^{\tld}\ .
\tag*{(1.2.1)}
\]
\textit{More generally (in appearance only!), if $\Phi_1$ is a subset of \textup{$\mathrm{Fl}(\mathcal{M})$} which is transfinitely left \textup{Q}-saturated, and such that}
\[
\Phi_0 \ \subset \ \Phi_1 \ \subset \ \Phi^{\tld} \ = \ \Phi_0{}^{\tld}\ ,
\]